% A146: the exact integral map and a positive-density reduction obstruction.
\subsection{Boundary reduction for infinitely supported elliptic classes}

\begin{statement}{A.146}{Proved}
{An infinite-primary lifting obstruction on the rank-two elliptic factor}
Retain the actual maps
\(D^\circ\longrightarrow E\xrightarrow{\phi}F\) of A.143.
Write
\[
 M_\ell=\varprojlim_j F(\Q)/\ell^jF(\Q),\qquad
                 \widehat{F(\Q)}=\prod_\ell M_\ell .
\tag{A.146.1}
\]
Let \(S\) be a set of odd coefficient primes for which
\[
       \#\{\ell\in S:\ell\le X\}=O(X^\delta)
                         \quad\text{for some }\delta<\tfrac12 .
\tag{A.146.2}
\]
Every non-torsion element of the closed subgroup
\[
 \mathcal M(S)=M_2\times
      \prod_{\substack{\ell\in S}}M_\ell\times
      \prod_{\substack{\ell\text{ odd}\\\ell\notin S}}\ell M_\ell
\tag{A.146.3}
\]
defines a proper elliptic Kummer section which has no full Selmer
lift to \(D^\circ\).

In particular, every non-torsion completed class supported on
such an \(S\cup\{2\}\) is excluded.  There are also excluded
classes with every primary component nonzero: for any
non-torsion \(R\in F(\Q)\), take
\[
             \eta=zR,\qquad z_2=1,\qquad z_\ell=\ell
                               \quad(\ell\text{ odd}).
\tag{A.146.4}
\]
Here \(z\) is the displayed element of
\(\widehat\Z=\prod_\ell\Z_\ell\).

The detecting primes are good primes satisfying
\[
 \begin{gathered}
 p\equiv3\pmod8,\qquad
       \ell\nmid p+1\quad(\ell\in S),\\
       \ell^2\nmid p+1\quad(\ell\text{ odd},\ \ell\notin S).
 \end{gathered}
\tag{A.146.5}
\]
Their density among all rational primes exists and equals
\[
 \mathfrak d(S)=\frac14
       \prod_{\ell\in S}\left(1-\frac1{\ell-1}\right)
       \prod_{\substack{\ell\text{ odd}\\\ell\notin S}}
                   \left(1-\frac1{\ell(\ell-1)}\right)>0 .
\tag{A.146.6}
\]
The local assertion used in the proof is reduction into the
boundary.  It does not require the local elliptic point itself
to be a torsion point or the local section itself to be cuspidal.
\end{statement}

\begin{proof}
\emph{The integral curve and its two boundary fibres.}
Choose a positive integer \(N\) divisible by every prime at
which one of \(2,k,q_0\) from A.143 is not a unit, and put
\(R_0=\Z[1/N]\).
The following projective plane quartic is a smooth proper model
of \(\overline D\) over \(R_0\):
\[
        aW^4=T(T-Z)(T-\tfrac12Z)Z
                     \quad\text{in }\mathbf P^2_{R_0}.
\tag{A.146.7}
\]
To verify smoothness, on \(Z\ne0\) a point with \(W\ne0\)
has nonzero \(W\)-derivative \(4aW^3\).
At \(W=0\), the three roots \(0,1,\tfrac12\) are distinct
and simple.  On \(Z=0\) there is just \([1:0:0]\), and the
derivative with respect to \(Z\) is a unit there.
These checks apply in every geometric fibre.
Its four boundary sections are
\[
 (T,W)=(0,0),(1,0),(\tfrac12,0),\qquad [1:0:0].
\tag{A.146.8}
\]

The Weierstrass equations of \(E\) and \(F\) are smooth
elliptic schemes over \(R_0\).
The formula
\[
                 x=n(2T-1),\qquad y=8k^2q_0W^2
\tag{A.146.9}
\]
defines the map to \(E\) on the affine part.
It also extends over the infinity section, as can be checked
without a specialization assumption.
On \(T=1\) in the projective coordinates of (A.146.7), write
\(w=W/T\), \(z_0=Z/T\).  Then
\[
             aw^4=z_0(1-z_0)(1-z_0/2),
\tag{A.146.10}
\]
so \(z_0/w^4\) is a unit near infinity.
There \(x\) has pole order four and \(y\) pole order six;
the projective Weierstrass coordinates \([x/y:1:1/y]\)
are regular and specialize to \([0:1:0]\).
This proves the extension of the literal map in (A.146.9).

The formula for \(\phi\) in (A.143.13) likewise extends over
these elliptic schemes.
Away from \((0,0)\) and the origin its denominators are units
on their usual affine charts.  Near \((0,0)\), the relation
\(y^2=x(x^2-n^2)\) makes \(x/y^2\) a unit;
the target chart \([u/v:1:1/v]\) is regular there and maps
that section to the origin.  Near the origin the same target
chart is regular, since \(x,y\) have pole orders two and three.
The induced differential satisfies
\[
                    \phi^*\frac{du}{2v}=2\frac{dx}{2y}.
\tag{A.146.11}
\]
It is therefore an etale degree-two isogeny in every fibre,
with the two distinct kernel sections \(O_E,(0,0)_E\).
The explicit inverse fibres in A.143 persist over every
algebraically closed residue field of characteristic not
dividing \(N\):
\[
 \begin{aligned}
 (\phi q)^{-1}(O_F)&=2\{T=\tfrac12\}+2\{T=\infty\},\\
 (\phi q)^{-1}((0,0)_F)&=2\{T=0\}+2\{T=1\}.
 \end{aligned}
\tag{A.146.12}
\]
Indeed, \(q^{-1}(E[2])\) consists exactly of (A.146.8),
by (A.146.9) and the infinity calculation.
At a finite boundary point \(W\) is a parameter and
\(y\) is a unit times \(W^2\), so the ramification index
of \(q\) is two.  At infinity the parameter \(x/y\) of
\(E\) has order two by (A.146.10).
The two fibre degrees are four, so there are no additional
geometric points in (A.146.12).

The actual map to the tripod is given over \(R_0\) by
\[
          t=-4T(T-1),\qquad
          [t:1]=[-4T(T-Z):Z^2].
\tag{A.146.13}
\]
The latter expression is a morphism on (A.146.7), including
its infinity point.  It sends the four sections (A.146.8)
respectively to \(0,0,1,\infty\).
Consequently a \(\Q_p\)-point of \(D^\circ\) whose image
on \(F\) reduces to \(O_F\) or \((0,0)_F\) reduces to a
boundary point of (A.146.7), and its tripod image is not in
\(U(\Z_p)\).  Properness supplies the reduction of every
such \(\Q_p\)-point.

\emph{The exact density.}
For each fixed cutoff \(Y\), impose (A.146.5) only for
odd primes \(\ell\le Y\).
The Chinese remainder theorem on reduced residue classes,
followed by Chebotarev for the resulting cyclotomic extension,
gives density
\[
 \mathfrak d_Y(S)=\frac14
       \prod_{\substack{\ell\le Y\\\ell\in S}}
                 \left(1-\frac1{\ell-1}\right)
       \prod_{\substack{\ell\le Y,\ \ell\text{ odd}\\\ell\notin S}}
                 \left(1-\frac1{\ell(\ell-1)}\right).
\tag{A.146.14}
\]
Each forbidden residue is \(-1\), modulo respectively
\(\ell\) or \(\ell^2\), and is a reduced residue.
The modulus eight condition contributes exactly \(1/4\).
Condition (A.146.2), by summation over dyadic intervals, implies
\(\sum_{\ell\in S}1/(\ell-1)<\infty\).
The corresponding sum for \(1/(\ell(\ell-1))\) converges
even over all positive integers \(\ell\ge3\).
Thus (A.146.14) converges to the positive product
(A.146.6).

Here are the uniform tail estimates that justify passage from
finite residue conditions to the full set of primes.
The exact Brun--Titchmarsh inequality of
Montgomery--Vaughan \cite[Theorem 2, (1.10), p.121]{LegMontgomeryVaughan1973}
gives, for \(q\le\sqrt X\),
\[
              \pi(X;q,-1)\le
                   \frac{4X}{\varphi(q)\log X}.
\tag{A.146.15}
\]
For example take the limit of their interval formula as its
left endpoint tends to zero.  Enlarging a harmless absolute
constant handles the bounded initial values of \(X\).
For the forbidden first powers at primes of \(S\), this gives
\[
 \#\{p\le X:\exists\ell\in S,\ Y<\ell\le\sqrt X,
                         \ \ell\mid p+1\}
 \le\frac{4X}{\log X}
                     \sum_{\substack{\ell\in S\\\ell>Y}}
                                 \frac1{\ell-1}.
\tag{A.146.16}
\]
For \(\sqrt X<\ell\le X+1\), count all integers in the
progression instead.  Each \(\ell\) allows at most
\(\sqrt X+2\) integers \(p\le X\), and (A.146.2)
bounds the total by
\[
                           O(X^{\delta+1/2}+X^\delta)
                                    =o(X/\log X).
\tag{A.146.17}
\]
Larger \(\ell\) cannot divide \(p+1\).

For the square conditions outside \(S\), apply (A.146.15)
with \(q=\ell^2\) when \(\ell\le X^{1/4}\).
Their tail is bounded by
\[
                 \frac{4X}{\log X}
                 \sum_{\substack{\ell>Y,\ \ell\text{ odd}\\\ell\notin S}}
                              \frac1{\ell(\ell-1)}.
\tag{A.146.18}
\]
For larger \(\ell\), count all possible integers rather than
primes, and even allow every integer divisor \(m>X^{1/4}\).
The bound is
\[
 \sum_{X^{1/4}<m\le\sqrt{X+1}}
                   \left(\frac{X+1}{m^2}+1\right)
                  =O(X^{3/4}+\sqrt X)=o(X/\log X).
\tag{A.146.19}
\]
Divide the four bounds by \(\pi(X)\sim X/\log X\),
first let \(X\) tend to infinity and then let \(Y\) tend
to infinity.  Both reciprocal tails vanish.
The upper and lower densities therefore coincide with
(A.146.6).  Natural density among primes implies the same
Dirichlet density, by partial summation.
Removing the finitely many bad model primes does not alter it.
For \(S=\varnothing\), the condition is simply that
\(p\equiv3\pmod8\) and \((p+1)/4\) is squarefree.

\emph{Localization and the reduction map.}
The surjection \(E(\Q)\to F(\Q)\) proved in A.143 induces
surjections on each primary completion.
For every odd \(\ell\notin S\), write
\(\eta_\ell=\ell\xi_\ell\), lift \(\xi_\ell\) to the
primary completion of \(E(\Q)\), and multiply that lift by
\(\ell\).  At the other primes choose any lift.
Taking their product gives
\[
 \widehat\phi(\theta)=\eta,\qquad
 \theta_\ell\in\ell\widehat{E(\Q)}_\ell
                    \quad(\ell\text{ odd},\ \ell\notin S).
\tag{A.146.20}
\]
No choice of a single rational multiplier is made here.

At every good prime \(p\equiv3\pmod8\), the odd-cubic
character-sum calculation of A.143 gives
\[
                \#E(\F_p)=\#F(\F_p)=p+1,
                   \qquad v_2(p+1)=2.
\tag{A.146.21}
\]
Reduction is a continuous homomorphism from the local compact
groups, hence also from the global profinite completions.
A homomorphism from a pro-\(\ell\) abelian group to a finite
group has image in its \(\ell\)-primary subgroup.
At a prime satisfying (A.146.5), every odd component in
\(S\) therefore reduces to zero; all other odd components
reduce to zero because their target Sylow subgroup has
exponent at most \(\ell\) and they are divisible by \(\ell\).
This includes the physical coefficient prime \(\ell=p\):
its reduction is zero because \(p\nmid p+1\), regardless
of its possibly nonzero formal-group component.

The two-primary subgroup of \(E(\F_p)\) has order four
and contains the four distinct reductions of \(E[2]\).
It is exactly their image.  Applying \(\phi\) to the
reduction of (A.146.20) consequently gives
\[
                 \operatorname{red}_p(\eta_p)
                       \in\{O_F,(0,0)_F\}
                                  \subset F(\F_p).
\tag{A.146.22}
\]
Only this reduction statement is needed.

\emph{The section obstruction.}
Suppose a full Selmer section on \(D^\circ\) has compact
elliptic Kummer class \(\eta\).
At a detecting prime, its local point or cusp maps to the
localization \(\eta_p\).  This follows from functoriality of
the Kummer maps and their local injectivity, as in A.143.
For an ordinary point, (A.146.12) and (A.146.22) force
reduction to a boundary point; (A.146.13) then makes its
tripod image nonintegral.  A local cusp already maps to a
tripod cusp and also supplies no integral ordinary point.
The pushed-forward full Selmer section on \(U\) thus fails
ordinary integrality at every prime in a set of positive
Dirichlet density.

The density-one integrality theorem AP-D.25, in the same
Selmer range used in A.143, forces this image section to be
globally cuspidal \cite[D.25]{LegAP2026}.
Intersect its cusp decomposition group with the actual open
subgroup for \(D^\circ\to U\).  As in A.143, the
intersection is a decomposition group at a boundary point
of \(\overline D\), and the section's surjection onto
\(G_\Q\) makes that boundary point rational.
Thus the original section is cuspidal on \(D^\circ\).
Its compact elliptic image is consequently \(O_F\) or
\((0,0)_F\), by (A.146.12).
Global Kummer injectivity on \(\widehat{F(\Q)}\), proved
in A.143, contradicts the non-torsion assumption on \(\eta\).

The actual proper Kummer sections for the displayed completion
classes are the continuous Tate-module cocycles described in
A.143.  For (A.146.4), a non-torsion rational point has a
nonzero free component over every \(\Z_\ell\); multiplication
by the nonzero scalar \(z_\ell\) preserves it.
This gives every-primary-component-nonzero examples within
\(\mathcal M(\varnothing)\).
Finally, infinite sets satisfying (A.146.2) are explicit:
start with \(\ell_1=3\) and recursively take the least prime
\(\ell_{j+1}>\ell_j^2\).
Then their counting function is \(O(\log\log X)\), by the
lower bound \(\ell_j\ge3^{2^{j-1}}\), and hence satisfies
(A.146.2).  Completed classes supported there give the
stated sparse infinite-support examples.
\end{proof}

The certificate \path{code/harmonic_infinite_primary_certificate.py}
checks the exact reduced-residue densities for finite mixed
conditions, the good special-fibre boundary maps, and the
complete finite-group reduction assertion at explicit primes.
Its finite counts do not replace the tail estimates
(A.146.16)--(A.146.19).
